\documentclass[11pt]{article}
\usepackage[margin=1.1in]{geometry}
\usepackage{amsmath,amssymb,url}
\newcommand{\grad}{\nabla}
\newcommand{\eps}{\varepsilon}
\setlength{\parskip}{4pt}

\title{\vspace{-2em}Constructive Solution of the Clay Navier--Stokes Problem}
\author{Claes Johnson\\[0.2em]
\small KTH Royal Institute of Technology, SE-100 44 Stockholm, Sweden\\
\small \texttt{claesjohnson@gmail.com}}
\date{\today\\[0.4em]\small\textit{written up with AI assistance (Claude, Anthropic)}}

\begin{document}
\maketitle

\begin{abstract}
The solution of alternative~(C) of the Clay Navier--Stokes problem announced in September 2026
represents, beyond its result, a new form of mathematics: one performed by AI as a digital
optimisation process, fitting a large set of network parameters to given data by forward and
backward propagation. Mathematics done this way constructs an object rather than asserting that one
exists, and AI in this sense may be seen as a form of automated thinking. We describe a far more
restricted version of the same thing, developed over a long period and implemented in the FEniCS
software: the automated solution of differential equations, ASDE, a self-learning computational
algorithm with dual-based output error control directed at minimal computational effort, whose
backward pass is the adjoint and is the same computation as backpropagation. We then present what
ASDE delivers when directed at alternative~(A): constructive existence of turbulent solutions of the
incompressible Navier--Stokes equations at high Reynolds number, for data given in advance, together
with qualitative uniqueness of their mean-value outputs, drawn quantitatively by computed stability
factors so that a functional can fail the test. We do not claim a discharge of (A) as that
alternative is formulated, since (A) asks for smoothness and the solutions that occur are not
smooth: their gradients grow like $\nu^{-1/2}$ while their velocities stay bounded, a
H\"older-$1/3$ field neither differentiable nor unbounded. We argue that this third case is what
the equations produce, and close with a scaling test which shows that the singular solution
constructed for (C) vanishes in the limit of small viscosity in which fluid mechanics lies: held to
a fixed admissible forcing class its amplitude, support and lifetime all tend to zero, and held at
fixed amplitude its forcing diverges like $\nu^{-1}$ and leaves the class.
\end{abstract}

\section{A new form of mathematics, and a restricted version of it}

On 8~September 2026 OpenAI presented a solution of alternative~(C) of the Navier--Stokes
Millennium problem~\cite{openai}: a proof that a singular solution exists, obtained with ten
thousand agents and of the order of $3\times10^{11}$ output tokens~\cite{buckmaster}.

An
AI is a digital optimisation process: a very large set of network parameters is fitted to given
data by forward and backward propagation, the backward pass carrying the sensitivity of an output
to everything that fed it. Mathematics done this way does not establish that an object exists and
leave it there; it \emph{constructs} one, by fitting. AI in this sense is a form of
\emph{automated thinking}, and the (C) solution is an instance: the singular velocity and pressure
fields are built first, to have the required growth, and the force is then defined as what is left
over --- the residual of the construction, arranged to be smooth.

Over a long period, together with students and colleagues, we have developed a far more restricted
version of the same thing: the \emph{automated solution of differential equations}, which we will
call ASDE, implemented in the FEniCS software~\cite{unicorn,massparallel,fenicshpc,fenicsmulti} and
now widely used. ASDE fits a discrete solution to a differential equation rather than a network to
data, and it does so by the same forward and backward passes. It is a small, exactly specified
corner of automated thinking, and it has one advantage over the general case: what it constructs
comes with a bound on the error in whatever is asked of it.

\textbf{What follows is a constructive solution directed to alternative~(A).} We do not claim it is a
solution of (A) as that alternative is formulated, and the difference should be stated at the outset
rather than discovered later. (A) asks, for \emph{all} admissible data and with $f\equiv0$, for a
solution that exists and is \emph{smooth} for all time. What ASDE delivers, for data given in
advance, is \emph{constructive existence} and \emph{qualitative uniqueness} of turbulent solutions.
It does not deliver smoothness, because the solutions that occur do not have it --- their gradients
grow like $\nu^{-1/2}$ while their velocities stay bounded (Section~\ref{sec:third}).

So this is offered as an answer to what (A) is \emph{about} --- whether the equations determine a
flow, and what of it they determine --- and not as a discharge of (A) as written. Whether the two
coincide is for the reader to judge, and it is the one judgement we would ask for.

None of the material is new. It was published in 2007, in Chapter~14 of a book written with Johan
Hoffman~\cite{cfd}, and it has been before the Institute before. What is new is the occasion.


\section{Basics of ASDE}\label{sec:asde}

ASDE is a self-learning computational algorithm with dual-based output error control, directed at
minimal computational effort. Its three parts are the following.

\paragraph{A discretisation that fits the equation.} Given a differential equation, a domain and
data, a finite element method produces a discrete solution $U$ whose \emph{residual} $R(U)$ measures
by how much it fails to satisfy the equation. Here the method is General Galerkin, G2: continuous
piecewise linear in space and time, $\mathrm{cG}(1)\mathrm{cG}(1)$, with least-squares control of
the residual~\cite{fehn}. Nothing is supplied by the user beyond equation, domain and data --- no
mesh chosen by hand, no turbulence model, no closure, no constant fitted to measurement. This
matters for what follows: a computation carrying a fitted turbulence model could not be used to say
what the Navier--Stokes equations contain, because the turbulence would have been put in by hand.

\paragraph{A backward pass that converts residual into error.} The error in an output $M(u)=\int_Q
u\cdot\psi$ is, to leading order, the residual tested against a \emph{dual weight function}
$\varphi$,
\begin{equation}\label{eq:dualrep}
M(u) - M(U) \;\approx\; \int_0^T\!\!\int_\Omega R(U)\cdot\varphi \,dx\,dt ,
\qquad
|M(u)-M(U)| \;\le\; S\,\|R(U)\|_{-1},
\end{equation}
with $S$ the stability factor. And $\varphi$ is obtained by solving the equation linearised about
the computed solution \emph{backward in time} from the output functional as data at $t=T$. This is
the same object, and the same computation, as the backward pass in the training of a neural
network: the adjoint of a forward evolution, propagating sensitivity from an output to everything
that fed it. ASDE and AI are doing the same linear algebra in opposite directions through the same
trajectory, and that is the sense in which ASDE is the restricted case.

\paragraph{Self-learning: adaptivity towards minimal effort.} The mesh is neither fixed in advance
nor refined uniformly. Refinement goes where the product of residual and dual weight is largest,
and stops when the resulting bound meets the tolerance asked. Adaptivity and error control are one
mechanism, not two: the computation that says how large the error is also says where to refine to
reduce it, and the result is the least computational work that achieves the stated accuracy. The
algorithm learns the mesh its own output requires.

\paragraph{What ASDE therefore produces.} Not an assertion that a solution exists, but an object:
the mesh with its cell count and smallest cell, the discretisation carrying the solution, the
residual measured in the norm the bound uses, a computed stability factor for each output, and for
each output a value with a bound on its error. That is the solution \emph{de facto}. The record of
the method and its computations is in~\cite{cfd}, with the residual-based dissipation
in~\cite{hjv}, the parameter-free formulation in~\cite{hjja}, the earlier statement
in~\cite{newapproach}, a case at aeronautical Reynolds number --- the DLR-F11 aircraft model,
time-resolved and adaptively refined --- in~\cite{dlrf11}, and two results the same computations
gave on the way: the resolution of d'Alembert's paradox~\cite{dalembert} and a theory of
flight~\cite{flight}. Further material is collected under the
label \emph{clay problem}~\cite{blog}.

\paragraph{Where it now stands.} The development since 2007 has gone to compressible flow and to
aeronautical Reynolds numbers, with parallel-in-time computation of turbulent
flow~\cite{paralleltime}, carried substantially by Johan Jansson, who has a leading role in
the AIAA CFD High Lift Prediction Workshop (HLPW-5 and HLPW-6), where computational predictions of
high-lift configurations are submitted blind and compared against wind-tunnel measurement. That is
the hardest ordinary test there is for a turbulence prediction: separated flow at aeronautical
Reynolds number, outputs that are exactly the mean-value functionals treated below, and a
comparison made before the measurement is known. ASDE is tested in the ordinary way, against
experiment, by people with no stake in this letter.


\section{ASDE applied to Navier--Stokes}\label{sec:ns}

The problem is the incompressible Navier--Stokes equations
\begin{equation}\label{eq:ns}
\dot u + (u\cdot\grad)u - \nu\Delta u + \grad p = f,\qquad \grad\cdot u = 0,
\end{equation}
on a domain $\Omega\subset\mathbb{R}^3$ with initial data $u_0$, forcing $f$ and boundary
conditions, at small viscosity, that is at high Reynolds number. The official statement, as set out
by Charles Fefferman~\cite{fefferman}, poses them instead on $\mathbb{R}^3$ or a periodic box with
the viscosity fixed, and asks whether solutions remain smooth for all time. Applied to it, ASDE delivers the
following two things.

\paragraph{Constructive existence of turbulent solutions.} A G2 computation of \eqref{eq:ns}
started from smooth potential data becomes turbulent and stays bounded. The stabilisation supplies
dissipation in proportion to the residual, so the dissipation is generated by the computation
rather than prescribed by a model; in the vanishing-viscosity limit the same scheme is an Euler
solver whose turbulent dissipation appears automatically. Existence here is not an assertion about
a solution set being non-empty: a solution is exhibited, for data given in advance, with the error
in each output bounded by \eqref{eq:dualrep}. Ask a non-constructive existence theorem for the drag
and there is nothing there to ask; ask a computed solution and it answers with a number and a bound.

There is more to this than error control. A computed solution is an elaborate explicit construction
of a velocity--pressure field --- millions of numbers on a spacetime mesh, every one available ---
and such an object can be taken apart: where vorticity concentrates and in what shape, how the
velocity gradient decomposes at a point, where energy passes from large scales to small, where the
dissipation actually sits as against where a model would put it. A solution concept whose objects
can be inspected generates conjectures; one whose objects cannot be exhibited does not. The triple
decomposition of the velocity gradient~\cite{triple}, and the energy stability analysis built on
it~\cite{energystab}, came out of looking closely at resolved velocity fields, and at measured
ones~\cite{musci}.

\paragraph{Qualitative uniqueness.} Classical uniqueness asks that the data determine the field
pointwise. At high Reynolds number that is false, and not marginally: two computations of the same
flow on independently generated meshes give velocity fields differing by order one at any chosen
point, and continuing to differ under refinement. What the data \emph{do} determine is the
qualitative content, and this is what the two computations are found to share: the drag agrees to a
few per cent~\cite{drag} and the lift likewise; the separation line sits in the same place; the wake
is organised the same way and sheds at the same frequency; the total dissipation rate is the same.

Qualitative uniqueness is the statement that this is so, with the line between the two lists drawn
\emph{quantitatively} rather than by assertion --- and that last clause is what keeps it from being
a soft version of uniqueness. The dual solution attaches to each functional a computed stability
factor: where it is moderate the output is determined by the data to a stated precision, and
refinement where the dual weight is large will improve it; where the factor is large the output is
not determined, and no refinement will make it so. The claim is falsifiable one functional at a
time, and a functional can fail it. That last case is not a failure of the method but an answer:
the equations do not determine that quantity, and no computation should be believed which claims
otherwise.

\paragraph{Turbulence is a necessity, not a possibility.} High Reynolds number flow starting out
laminar develops shear instabilities with sharp velocity gradients, which are eventually damped by
viscosity in turbulent flow. All high Reynolds number laminar flow is thus unstable and cannot
continue to exist over time; it turns turbulent in order to continue to exist. What ASDE constructs
is therefore not one possibility among several but what the equations must do.


\section{Why smoothness is not what the solutions have}\label{sec:third}

Since (A) as written asks for smoothness, it is worth being precise about what the computed
solutions are instead, because they are neither smooth nor singular. Dissipation at a rate
independent of the viscosity gives gradients $|\grad u|\sim\nu^{-1/2}$, and the scales at which
that gradient is realised follow from the same constant:
\begin{equation}\label{eq:delta}
\eta \sim \Big(\frac{\nu^3}{\eps}\Big)^{1/4} \sim \nu^{3/4},
\qquad
\delta u(\eta) \sim (\eps\,\eta)^{1/3} \sim \nu^{1/4},
\qquad
\frac{\delta u(\eta)}{\eta} \sim \nu^{-1/2},
\end{equation}
the first two being the Kolmogorov--Obukhov scaling~\cite{kolmogorov41,obukhov41}, refined by the
intermittency corrections of~\cite{kolmogorov62,obukhov62} and the multifractal picture
of~\cite{frischparisi}, and the middle relation being H\"older continuity with exponent $1/3$, the
Onsager exponent. As $\nu\to0$ the structures refine without limit while the velocity stays bounded
at every scale: the gradient diverges, so the solution is not smooth, and the velocity does not, so
it is not singular. A H\"older-$1/3$ field is exactly the object that is neither differentiable nor
unbounded, which is why no amount of regularity theory reaches it.

This category is not conjectural. By the Onsager theorem~\cite{onsager,isett}, proved after the
partial results of~\cite{eyink,cet,dls,bdlis}, non-smoothness is the \emph{mechanism} by which a
turbulent flow dissipates energy at a rate that does not vanish with the viscosity --- the
dissipation anomaly, observed since Taylor~\cite{taylor} and still under active
examination~\cite{iyer} --- and a local energy equation can be written in which the non-smoothness
is the direct cause of the dissipation~\cite{duchonrobert}. Existence of weak solutions has been
available since Leray~\cite{leray,lions}, and for small data, the low Reynolds number case for
suitably scaled problems, global existence of a smooth solution has been known for sixty
years~\cite{fujitakato,kato}. The difficulty lies at the other end, and that is the end ASDE
addresses.


\section{The (C) solution under the viscosity test}\label{sec:nu}

Two observations about the announced solution, offered without disputing it.

\paragraph{(C) is not the negation of (A).} Fefferman's four statements do not form two pairs of
negations. (A) and (B) are posed with $f\equiv0$ and quantify over \emph{all} admissible data;
(C) and (D) ask for the \emph{existence} of admissible data \emph{together with an admissible
smooth forcing} $f$ for which smoothness fails. So a proof of (C) leaves (A) untouched: it does not
show that a fluid tears itself into a singularity, only that something can be done to it that does.
The paper states its own scope in just these terms --- it establishes (C), with the periodic case
giving (D) --- and does not claim (A) or (B). So (A) remains open after it, which is why a solution
of (A) is worth presenting.

\paragraph{The viscosity test.} The test is general and applies to any construction exhibiting a
singularity at prescribed forcing. Navier--Stokes rescales: if $v$ solves the equations at viscosity
one with forcing $g$, then
\begin{equation}\label{eq:rescale}
u(x,t) = a\,v(bx,ct),\qquad a = \nu b,\quad c = ab,\quad f = a^2 b\,g,
\end{equation}
solves them at viscosity $\nu$, and the map is onto: \emph{one} unit-viscosity blowup generates the
entire family. A theorem asserting blowup ``for every $\nu>0$'' is therefore not a family of
results. It is one result, rescaled, containing exactly what its $\nu=1$ member contains.

Everything then turns on what \eqref{eq:rescale} does to the constants, and since it is the forcing
class that the official statement fixes, two normalisations bear on it. Hold the solution at size
one, $a=1$: then the forcing amplitude is $a^2b=\nu^{-1}$, so the force required \emph{blows up} as
the viscosity falls and leaves any fixed admissible class. Or hold the forcing in a fixed class,
$a^2b=1$: then $b=\nu^{-2/3}$ and
\begin{equation}\label{eq:shrink}
\text{amplitude } a = \nu^{1/3}\to 0,\qquad
\text{length } b^{-1} = \nu^{2/3}\to 0,\qquad
\text{time } c^{-1} = \nu^{1/3}\to 0 .
\end{equation}
The singularity shrinks to nothing: at $\nu=10^{-6}$ the velocity reaches $10^{-2}$ on a scale of
$10^{-4}$ before blowing up. Neither branch survives the limit --- one exits the admissible data,
the other converges to the zero solution --- and there is no third option. Set this beside the flow
the problem is about: turbulence at small viscosity has amplitude \emph{fixed} and structure
refining, by \eqref{eq:delta}. The constructed family has amplitude vanishing like $\nu^{1/3}$,
which is the opposite behaviour, or else forcing diverging like $\nu^{-1}$, which is not a fluid at
all. The independent exclusion of the mechanism under analytic forcing~\cite{civ} points the same
way from a different direction.


\section{In closing}

Two constructions, then, of the same formal kind and opposite in direction. The (C) solution fixes
the singularity wanted and solves backwards for a force that produces it, the force being whatever
the residual requires. ASDE fixes the data a flow actually has and solves forwards, the residual
becoming the bound on what may be asked of the result. Both are automated, both exhibit, and only
the second exhibits a fluid.

We do not ask for the problem to be restated, and we do not claim to have discharged (A) as it is
formulated. What we present is what ASDE delivers when directed at it --- constructive existence of
turbulent solutions and qualitative uniqueness of their mean-value outputs, for specific data, with
error bounds the computation itself produces --- and we leave the reader to judge how near that comes
to the question the problem was meant to ask. The material has been in print
for some twenty years, from~\cite{newapproach} onward and set out in Chapter~14 of~\cite{cfd}; the
programme has continued since into compressible flow and aeronautical Reynolds numbers. What has changed is that a solution of (C) now exists, by methods
that are themselves a new form of mathematics, and that makes it worth asking what the same kind of
mathematics delivers on (A).

\begin{thebibliography}{9}
\bibitem{fujitakato} H.~Fujita and T.~Kato, \emph{On the Navier--Stokes initial value problem~I},
Arch.\ Rational Mech.\ Anal.\ \textbf{16} (1964) 269--315.

\bibitem{kato} T.~Kato, \emph{Strong $L^p$-solutions of the Navier--Stokes equation in
$\mathbb{R}^m$, with applications to weak solutions}, Math.\ Z.\ \textbf{187} (1984) 471--480.

\bibitem{kolmogorov41} A.~N.~Kolmogorov, \emph{The local structure of turbulence in incompressible
viscous fluid for very large Reynolds numbers}, Dokl.\ Akad.\ Nauk SSSR \textbf{30} (1941) 301--305.

\bibitem{obukhov41} A.~M.~Obukhov, \emph{On the distribution of energy in the spectrum of turbulent
flow}, Bull.\ Acad.\ Sci.\ USSR, Geog.\ Geophys.\ Ser.\ \textbf{5} (1941) 453--466.

\bibitem{kolmogorov62} A.~N.~Kolmogorov, \emph{A refinement of previous hypotheses concerning the
local structure of turbulence in a viscous incompressible fluid at high Reynolds number}, J.\ Fluid
Mech.\ \textbf{13} (1962) 82--85.

\bibitem{obukhov62} A.~M.~Obukhov, \emph{Some specific features of atmospheric turbulence}, J.\
Geophys.\ Res.\ \textbf{67} (1962) 3011--3014.

\bibitem{frischparisi} U.~Frisch and G.~Parisi, \emph{On the singularity structure of fully developed
turbulence}, in M.~Ghil, R.~Benzi and G.~Parisi (eds.), \emph{Turbulence and Predictability in
Geophysical Fluid Dynamics and Climate Dynamics}, North-Holland, 1985.

\bibitem{musci} B.~Musci, B.~Dubrulle, J.~LeBris, D.~Geneste, P.~Braganca, J.-M.~Foucaut, C.~Cuvier
and A.~Cheminet, \emph{Experimental measurement of the vorticity--strain alignment around extreme
energy transfer events}, J.\ Fluid Mech.\ \textbf{1016} (2025) A53.

\bibitem{iyer} K.~P.~Iyer, T.~D.~Drivas, G.~L.~Eyink and K.~R.~Sreenivasan, \emph{Turbulence without
walls: whither the zeroth law of turbulence?}, Phys.\ Rev.\ Lett.\ \textbf{135} (2025) 134001.

\bibitem{taylor} G.~I.~Taylor, \emph{Statistical theory of turbulence}, Proc.\ R.\ Soc.\ Lond.\ A
\textbf{151} (1935) 421--444.

\bibitem{eyink} G.~L.~Eyink, \emph{Energy dissipation without viscosity in ideal hydrodynamics~I.
Fourier analysis and local energy transfer}, Physica~D \textbf{78} (1994) 222--240.

\bibitem{dls} C.~De~Lellis and L.~Sz\'ekelyhidi~Jr., \emph{Dissipative Euler flows and Onsager's
conjecture}, J.\ Eur.\ Math.\ Soc.\ \textbf{16} (2014) 1467--1505.

\bibitem{bdlis} T.~Buckmaster, C.~De~Lellis, P.~Isett and L.~Sz\'ekelyhidi~Jr., \emph{Anomalous
dissipation for $1/5$-H\"older Euler flows}, Ann.\ of Math.\ \textbf{182} (2015) 127--172.

\bibitem{duchonrobert} J.~Duchon and R.~Robert, \emph{Inertial energy dissipation for weak solutions
of incompressible Euler and Navier--Stokes equations}, Nonlinearity \textbf{13} (2000) 249--255.

\bibitem{fehn} N.~Fehn, M.~Kronbichler and G.~Lube, \emph{From anomalous dissipation through Euler
singularities to stabilized finite element methods for turbulent flows}, Flow Turbul.\ Combust.\
\textbf{115} (2025) 347--388.

\bibitem{triple} J.~Kronborg and J.~Hoffman, \emph{The triple decomposition of the velocity gradient
tensor as a standardized real Schur form}, Phys.\ Fluids \textbf{35} (2023) 031703.

\bibitem{energystab} J.~Hoffman, \emph{Energy stability analysis of turbulent incompressible flow based
on the triple decomposition of the velocity gradient tensor}, Phys.\ Fluids \textbf{33} (2021) 081707.

\bibitem{hjv} J.~Hoffman, J.~Jansson and R.~Vilela de Abreu, \emph{Adaptive modeling of turbulent
flow with residual based turbulent kinetic energy dissipation}, Comput.\ Methods Appl.\ Mech.\
Engrg.\ \textbf{200} (2011) 2758--2767.

\bibitem{hjja} J.~Hoffman, J.~Jansson, N.~Jansson and R.~Vilela de Abreu, \emph{Towards a
parameter-free method for high Reynolds number turbulent flow simulation based on adaptive finite
element approximation}, Comput.\ Methods Appl.\ Mech.\ Engrg.\ \textbf{288} (2015) 60--74.

\bibitem{unicorn} J.~Hoffman, J.~Jansson \emph{et al.}, \emph{Unicorn: Parallel adaptive finite
  element simulation of turbulent flow and fluid--structure interaction}, Comput.\ Fluids, 2013.
\bibitem{massparallel} N.~Jansson, J.~Hoffman and J.~Jansson, \emph{Framework for massively parallel
  adaptive finite element computational fluid dynamics}, SIAM J.\ Sci.\ Comput., 2012.
\bibitem{fenicshpc} J.~Hoffman, J.~Jansson and N.~Jansson, \emph{FEniCS-HPC: Automated predictive
  high-performance finite element computing with aerodynamics applications}, Lecture Notes in
  Computer Science, Springer, 2016.
\bibitem{fenicsmulti} J.~Hoffman, J.~Jansson \emph{et al.}, \emph{FEniCS-HPC: Coupled multiphysics
  in computational fluid dynamics}, Lecture Notes in Computer Science, Springer, 2017.
\bibitem{paralleltime} J.~Jansson and J.~Hoffman, \emph{Direct FEM parallel-in-time computation of
  turbulent flow}, 2017.
\bibitem{dlrf11} J.~Jansson, J.~Hoffman, N.~Jansson and R.~Vilela de Abreu, \emph{Time-resolved
adaptive FEM simulation of the DLR-F11 aircraft model at high Reynolds number}, 2014.

\bibitem{newapproach} J.~Hoffman and C.~Johnson, \emph{A new approach to computational turbulence
modeling}, Comput.\ Methods Appl.\ Mech.\ Engrg.\ \textbf{195} (2006) 2865--2880.

\bibitem{openai} OpenAI, \emph{Finite time blowup for Navier--Stokes}, preprint, 2026.
\newblock \url{https://cdn.openai.com/pdf/32d9f210-8b73-45e0-91bc-82a30aef8a9a/navier-stokes.pdf}

\bibitem{buckmaster} T.~Buckmaster, \emph{Statement}, 2026.
\newblock \url{https://cims.nyu.edu/~tristanb/statement.pdf}

\bibitem{civ} P.~Constantin, M.~Ignatova and V.~Vicol, \emph{Regularity of asymptotically
axisymmetric solutions to the 3D Navier--Stokes equations with analytic forcing},
arXiv:2609.20803, 2026.

\bibitem{onsager} L.~Onsager, \emph{Statistical hydrodynamics}, Nuovo Cimento \textbf{6},
Suppl.~2 (1949) 279--287.

\bibitem{fefferman} C.~Fefferman, \emph{Existence and smoothness of the Navier--Stokes equation},
Official Problem Description, Clay Mathematics Institute, 2000.
\bibitem{leray} J.~Leray, \emph{Sur le mouvement d'un liquide visqueux emplissant l'espace},
Acta Math.\ \textbf{63} (1934) 193--248.
\bibitem{cet} P.~Constantin, W.~E and E.~S.~Titi, \emph{Onsager's conjecture on the energy
conservation for solutions of Euler's equation}, Comm.\ Math.\ Phys.\ \textbf{165} (1994) 207--209.
\bibitem{isett} P.~Isett, \emph{A proof of Onsager's conjecture}, Ann.\ of Math.\ \textbf{188}
(2018) 871--963.
\bibitem{cfd} J.~Hoffman and C.~Johnson, \emph{Computational Turbulent Incompressible Flow},
Applied Mathematics: Body and Soul, vol.~4, Springer, 2007.
\bibitem{dalembert} J.~Hoffman and C.~Johnson, \emph{Resolution of d'Alembert's paradox},
J.\ Math.\ Fluid Mech.\ \textbf{12} (2010) 321--334.
\bibitem{flight} J.~Hoffman, J.~Jansson and C.~Johnson, \emph{New theory of flight},
J.\ Math.\ Fluid Mech.\ \textbf{18} (2016) 219--241.
\bibitem{drag} J.~Hoffman, \emph{Efficient computation of mean drag for the subcritical flow past a
circular cylinder using general Galerkin G2}, Int.\ J.\ Numer.\ Meth.\ Fluids \textbf{59} (2009)
1241--1258.
\bibitem{lions} J.-L.~Lions, \emph{Quelques m\'ethodes de r\'esolution des probl\`emes aux
limites non lin\'eaires}, Dunod, Paris, 1969.

\bibitem{blog} C.~Johnson, posts under the label \emph{clay problem},
\url{https://claesjohnson.blogspot.com/search/label/clay\%20problem}.
\end{thebibliography}
\end{document}
